\section{Z-Chain: the identity rollup}
\label{sec:zchain-rollup}

\subsection{Why a rollup}

Sections~\ref{sec:tx-auth-envelope} and \ref{sec:pq-permit} both leave one
question open: \emph{where does the verifier get the ML-DSA public key}?
\texttt{TxAuthEnvelope.PublicKeyRef} is a 32-byte hash and
\texttt{PQPermit.OwnerAccountID} is a 48-byte identifier; neither is the
1{,}952-byte (or 60.6-KB expanded) pubkey the lattice verifier needs.

The naive answer --- carry the pubkey on the wire with every transaction or
permit --- is unaffordable. ML-DSA-65 pubkeys are 1{,}952 B; an
expanded pubkey is 60.6 KB. A chain with $10^8$ daily transactions at
1{,}952 B per pubkey is moving $186$ GB of pubkey duplication \emph{per day}.

The right answer is an identity rollup: pubkeys are registered once,
revoked or rotated as needed, and committed via a Merkle root that
\texttt{TxAuthEnvelope.ZIdentityRoot} binds. The verifier looks up the
pubkey by \texttt{PublicKeyRef} against the rollup state at the bound
root, with a small inclusion proof.

Z-Chain (\texttt{luxfi/consensus@v1.23.7},
\texttt{protocol/zchain/registry/}) is that rollup. HIP-0078 §4 defines its
shape.

\subsection{PQKeyRecord}

The leaf of the identity rollup is \texttt{PQKeyRecord}
(\texttt{protocol/zchain/registry/key\_record.go}):

\begin{lstlisting}[language=Go,caption={PQKeyRecord leaf. Source:
\texttt{protocol/zchain/registry/key\_record.go}.}]
type PQKeyRecord struct {
    AccountID      [48]byte    // owner identifier
    SchemeID       SigSchemeID // ML-DSA-65 etc.
    CompactPubkey  []byte      // seed form (1952 B for ML-DSA-65)
    ExpandedRef    [32]byte    // SHA3-256 of the expanded-pubkey form
    Status         KeyStatus   // Active / Revoked / Rotated
    EpochRegistered uint64
    EpochExpiry     uint64     // 0 = no expiry
}
\end{lstlisting}

\texttt{Status} is one of \texttt{KeyStatusActive}, \texttt{KeyStatusRevoked},
\texttt{KeyStatusRotated}. The verifier accepts only \texttt{KeyStatusActive}
records. \texttt{ExpandedRef} is the link into the side-channel pubkey store
(Section~\ref{sec:pq-permit}); the precompile dereferences this when it
needs the expanded form.

The leaf hash \texttt{PQKeyRecord.Hash()} is TupleHash256 over the leaf
fields under customization \texttt{LUX\_PQ\_KEY\_RECORD\_V1}. Two leaves
differing in any field have different hashes; therefore the Merkle tree
detects any modification.

\subsection{The five operations}

The Z-Chain accepts five operation types
(\texttt{protocol/zchain/registry/ops.go}). Each carries a TupleHash256
signing digest under an op-specific customization, plus a FIPS 204 §5.2
\texttt{ctx} of \texttt{LUX\_ZCHAIN\_REGISTRY\_V1} (\texttt{ops.go:43}) on
every ML-DSA signature. The op-specific customization plus the registry-wide
\texttt{ctx} are independent layers of domain separation; replaying a
signature from one op as another fails on both.

\begin{table}[H]
\centering
\caption{Z-Chain registry operations. Source:
\texttt{protocol/zchain/registry/ops.go}.}
\label{tab:zchain-ops}
\begin{tabular}{llp{7cm}}
\toprule
Op & Customization & Action \\
\midrule
\texttt{OpRegisterKey}      & \texttt{LUX\_OP\_REGISTER\_KEY\_V1}      & Insert fresh key; proof of possession self-sig. \\
\texttt{OpRotateKey}        & \texttt{LUX\_OP\_ROTATE\_KEY\_V1}        & Replace one key with another; sig under outgoing key. \\
\texttt{OpRevokeKey}        & \texttt{LUX\_OP\_REVOKE\_KEY\_V1}        & Mark a key revoked; sig under the revoked key (plus profile-pinned reason). \\
\texttt{OpAuthorizeSession} & \texttt{LUX\_OP\_AUTHORIZE\_SESSION\_V1} & Bind an ephemeral session pubkey to a long-term account. \\
\texttt{OpCommitTxAuthBatch} & \texttt{LUX\_OP\_COMMIT\_TX\_AUTH\_BATCH\_V1} & Commit a batch of \texttt{TxAuthEnvelope} hashes for later inclusion proofs. \\
\texttt{OpCommitPermitBatch} & \texttt{LUX\_OP\_COMMIT\_PERMIT\_BATCH\_V1}   & Commit a batch of \texttt{PQPermit} hashes for later inclusion proofs. \\
\bottomrule
\end{tabular}
\end{table}

Strictly six operations, not five, once \texttt{OpCommitPermitBatch} is
counted alongside \texttt{OpCommitTxAuthBatch}. The HIP-0078 informal
``five-op'' framing groups the two commit ops because they share their
verifier; the canonical six are documented in \texttt{ops.go}. We carry the
informal framing in the paper title but use the precise count here.

\subsubsection{OpRegisterKey}

The proof of possession is the canonical: a registrant signs their own
\texttt{AccountID} transcript under the key they're registering. Without
this, an attacker could re-register a victim's pubkey as their own and
intercept the victim's future transactions. The signing digest
(\texttt{ops.go:66}) binds:

\[
\text{digest} = \texttt{TupleHash256}\left(\begin{aligned}
& \texttt{"Lux/OpRegisterKey/v1"}, \\
& \texttt{NewRecord.AccountID}, \\
& \texttt{NewRecord.Hash()}
\end{aligned}\right) \quad \text{cust = LUX\_OP\_REGISTER\_KEY\_V1}.
\]

The customization plus the record hash plus the AccountID bind every byte
that defines ``this key, for this account, under this op''.

\subsubsection{OpRotateKey}

\texttt{OpRotateKey} binds the outgoing record's hash, the incoming record's
hash, and the AccountID. The signature is under the \emph{outgoing} key:
the existing key authorises its own succession. The verifier checks the
chain (\texttt{outgoing.AccountID == incoming.AccountID},
\texttt{outgoing.SchemeID == incoming.SchemeID}, \texttt{outgoing.Status ==
Active}, \texttt{incoming.Status == Active}), then runs ML-DSA verify against
the outgoing key. Atomically, the outgoing record transitions to
\texttt{KeyStatusRotated}, and the incoming record becomes the active key
for the AccountID.

\subsubsection{OpRevokeKey}

\texttt{OpRevokeKey} carries a profile-pinned \texttt{Reason} byte
(compromise, lost, lifecycle, governance). The signature is under the
revoked key. After acceptance, the record's \texttt{Status} is
\texttt{KeyStatusRevoked} and the \texttt{RevocationRoot} aggregator (below)
includes the leaf.

\subsubsection{OpAuthorizeSession}

A session key is an ephemeral pubkey with a TTL and a capability scope. The
holder of the long-term identity key signs an authorisation binding the
ephemeral pubkey hash, the expiry epoch, and the capability bitmap. The
session key can sign transactions whose
\texttt{TxAuthEnvelope.WalletSchemeID} matches the session scheme; the
verifier looks up the session key against \texttt{SessionKeyRoot} and
confirms the binding. Session keys do not appear in \texttt{IdentityRoot};
they appear in \texttt{SessionKeyRoot} only. The two roots are independent
to support cheap session-key rotation without changing
\texttt{IdentityRoot}.

\subsubsection{OpCommitTxAuthBatch and OpCommitPermitBatch}

These two ops commit Merkle roots over batches of accepted
\texttt{TxAuthEnvelope} or \texttt{PQPermit} hashes, so the execution-side
verifier on later chains (e.g., the EVM verifying a permit inside a
transaction) can supply a Merkle inclusion proof rather than the full
envelope. The batch root is signed by the consensus committee under the
\texttt{FinalitySchemeID}; the committee is the same one running the
Pulsar finality lane (Section~\ref{sec:qchain-finality}).

\subsection{ZRegistryRoots}

The seven-root aggregator (\texttt{registry/roots.go:40}):

\begin{lstlisting}[language=Go,caption={ZRegistryRoots aggregator. Source:
\texttt{protocol/zchain/registry/roots.go:40}.}]
type ZRegistryRoots struct {
    IdentityRoot       [48]byte   // active PQKeyRecord set
    AccountKeyRoot     [48]byte   // same set indexed by AccountID
    RevocationRoot     [48]byte   // revoked / rotated AccountIDs
    SessionKeyRoot     [48]byte   // authorised session keys
    ContractPolicyRoot [48]byte   // per-contract policy hashes
    TxAuthRoot         [48]byte   // OpCommitTxAuthBatch roots
    PermitAuthRoot     [48]byte   // OpCommitPermitBatch roots
}

func (r *ZRegistryRoots) Hash() [48]byte {
    parts := [][]byte{
        []byte("Lux/ZRegistryRoots/v1"),
        r.IdentityRoot[:],  r.AccountKeyRoot[:],
        r.RevocationRoot[:], r.SessionKeyRoot[:],
        r.ContractPolicyRoot[:], r.TxAuthRoot[:],
        r.PermitAuthRoot[:],
    }
    return tupleHash48(parts, "LUX_ZCHAIN_REGISTRY_ROOTS_V1")
}
\end{lstlisting}

The aggregator hash rides as \texttt{EpochCommitment.zchain\_state\_root} on
the Q-Chain. Q-Chain does not look inside; it consumes the 48-byte digest
and binds it into \texttt{ComputeRoundDigest}
(Section~\ref{sec:qchain-finality}). Adding a root bumps the customization
\texttt{V1 $\to$ V2} and is a hard fork --- the wire layout and the bound
set are both pinned by the tag.

The seven roots are decomposed because the verifier hot path consumes only
a subset:

\begin{itemize}[leftmargin=2em,nosep]
  \item \texttt{IdentityRoot} is consumed by \texttt{TxAuthEnvelope}
    verification (Step 6 in Section~\ref{sec:tx-auth-envelope}).
  \item \texttt{AccountKeyRoot} is the AccountID-indexed lookup tree for
    execution-side calls.
  \item \texttt{RevocationRoot} is consumed when checking that an
    AccountID has no live revocation entries.
  \item \texttt{SessionKeyRoot} is consumed when validating a
    session-key-signed envelope.
  \item \texttt{ContractPolicyRoot} is consumed by
    \texttt{ContractAuthProfile} lookups.
  \item \texttt{TxAuthRoot} and \texttt{PermitAuthRoot} support
    inclusion-proof verification of historical envelopes.
\end{itemize}

\subsection{VerifyAuthPassed: the hot-path inclusion verifier}

\texttt{VerifyAuthPassed} (\texttt{registry/auth\_proof.go:118}) is the
verifier the EVM (and S-VM / X-VM) calls when it needs to confirm that a
\texttt{TxAuthEnvelope} or \texttt{PQPermit} hash was committed by the
Z-Chain. The function signature is roughly:

\begin{lstlisting}[language=Go,caption={\texttt{VerifyAuthPassed} signature.
Source: \texttt{protocol/zchain/registry/auth\_proof.go:118}.}]
func VerifyAuthPassed(
    profile  *config.ChainSecurityProfile,
    root     [48]byte,           // ZRegistryRoots.TxAuthRoot or PermitAuthRoot
    leaf     [48]byte,           // hash of the envelope / permit
    proof    MerkleProof,
) (bool, error)
\end{lstlisting}

The verifier:

\begin{enumerate}[leftmargin=2em,nosep]
  \item Confirms the proof's \texttt{HashSuiteID} matches
    \texttt{profile.HashSuiteID}.
  \item Walks the proof, reconstructing the root via the leaf cust
    \texttt{LUX\_MERKLE\_LEAF\_V1} and the node cust
    \texttt{LUX\_MERKLE\_NODE\_V1}.
  \item Returns \texttt{(true, nil)} iff the reconstructed root equals the
    bound root.
\end{enumerate}

The two customizations are critical: a leaf hash cannot be confused with an
internal-node hash, so a proof cannot be smuggled by re-interpreting a leaf
as an internal node (the classical ``second preimage on Merkle trees''
attack). cSHAKE256 with distinct customizations prevents the collision
structurally.

\subsection{Why this is a rollup, not a state tree}

The seven roots aggregate \emph{commitments}, not state. The full identity
state lives in the Z-Chain's own state tree. The seven roots are
EpochCommitment-level digests that the Q-Chain finalises. This is the same
pattern an optimistic / ZK rollup follows on Ethereum: the rollup chain runs
its own VM and produces a commitment that the L1 chain accepts.

The difference is that Z-Chain is co-validated by the same committee that
validates the C-Chain, P-Chain, and X-Chain --- there is no separate rollup
operator. The finality lane (Section~\ref{sec:qchain-finality}) signs every
Z-Chain commitment with Pulsar-65, the same lane that finalises the other
chains. The result is a single PQ identity layer shared by every execution
chain on Lux.
